\documentclass{article}
\usepackage{graphicx} 
\usepackage{float}

\title{DataMining Report}
\author{Deen Willemsen, Laurens van Ballegooijen }
\date{October 2026}

\begin{document}
\maketitle

\begin{abstract}
Large language models can now write hotel reviews that are hard to tell apart
from reviews written by real guests, which makes AI-generated opinion spam a
growing problem for review platforms. In this study we test whether
traditional machine learning methods can detect such reviews, using the 2{,}000
English reviews of the MAiDE-up dataset: 1{,}000 genuine Booking.com reviews and
1{,}000 reviews generated by GPT-4. The reviews are represented as a
bag-of-words, and we compare five classifiers: multinomial naive Bayes, Lasso
logistic regression, a single classification tree, a random forest and gradient
boosting. Hyperparameters were tuned on 75\% of the data, and the final models
were evaluated on a held-out test set of 500 reviews. Naive Bayes reached the
highest test accuracy (0.918), but it did not differ significantly from logistic
regression, the random forest (0.914) or gradient boosting (0.908) according to
McNemar's test. All four were significantly more accurate than the single tree
(0.794). The random forest was not significantly better on negative than on
positive reviews. The most predictive terms suggest that AI-generated reviews
use generic, polished phrases such as ``left much to be desired'' and ``in the
heart of the city'', while genuine reviews describe concrete events, such as
what was booked or what the staff said.
\end{abstract}

\section{Introduction}
Online hotel reviews play an important role in helping consumers make decisions about hotels and other services. However, the increasing ability of large language models (LLMs) to generate fluent and human-like text has created new opportunities for producing deceptive reviews. Distinguishing between genuine reviews written by humans and reviews generated by an LLM is therefore an increasingly relevant classification problem.

Ignat et al. (2025) investigate this problem in their work, \textit{MAiDE-up: Multilingual Deception Detection of AI-generated Hotel Reviews}. The authors introduce the MAiDE-up dataset, which contains 10,000 real human-written hotel reviews and 10,000 AI-generated hotel reviews across ten languages. The dataset was created to investigate linguistic differences between human-written and LLM-generated hotel reviews and to evaluate methods for detecting AI-generated reviews. \cite{ignat-etal-2025-maide}

In this study, we focus on the English-language subset of the MAiDE-up dataset. We investigate whether traditional machine learning methods can distinguish between genuine and AI-generated hotel reviews. The review texts are represented using a bag-of-words representation, after which five different classification approaches are evaluated: Multinomial Naive Bayes, logistic regression with a Lasso penalty, a single classification tree, random forests, and gradient boosting.

The aim of this study is to compare the performance of these classifiers in detecting AI-generated hotel reviews and to investigate which textual features are most useful for distinguishing between human-written and AI-generated reviews.

\section{Data}
The data used in this study comes from the MAiDE-up (Multilingual Ai-generateD fakE reviews) dataset introduced by Ignat et al. (2025). The dataset was created to study the differences between human-written and LLM-generated hotel reviews. It contains 20,000 hotel reviews, of which 10,000 are genuine human-written reviews and 10,000 are AI-generated reviews. The dataset is balanced across language, hotel location, and sentiment. The ten languages represented in the dataset are Chinese, English, French, German, Italian, Korean, Romanian, Russian, Spanish, and Turkish. \cite{ignat-etal-2025-maide}

The genuine reviews were collected from Booking.com. The authors collected reviews from hotels in ten cities: Ankara, Beijing, Berlin, Bucharest, Madrid, New Delhi, Paris, Rome, Seoul, and Washington. To obtain a balanced dataset, they identified 250 hotels in each city and collected up to 50 reviews per hotel. The reviews were also balanced with respect to sentiment. Booking.com scores from 1 to 10 were used, where scores from 1 to 6 were considered negative and scores from 7 to 10 positive. For each language, the dataset contains 500 positive and 500 negative reviews. The authors additionally performed automatic and manual quality checks on the collected reviews. \cite{ignat-etal-2025-maide}

The AI-generated reviews were created using GPT-4. The generated reviews were constructed to follow the same distributions of language, location, and sentiment as the genuine reviews. The generation process used prompts specifying the language, hotel location, and sentiment of the review. The prompts instructed the model to generate a review consisting of an ``upside'' and a ``downside'' section, together with an overall review score. The authors also used few-shot prompting and included real hotel reviews in the generation process. Quality checks were performed by native speakers, who evaluated generated reviews for readability, syntax, and style. \cite{ignat-etal-2025-maide}

For this study, we restrict the dataset to the English-language reviews. This results in 2,000 reviews: 1,000 genuine human-written reviews and 1,000 AI-generated reviews. The resulting dataset is therefore balanced with respect to the target class. The reviews contain two text fields, an ``upside'' and a ``downside'', which correspond to the two parts of the review format used by Booking.com and the generation procedure described by Ignat et al. (2025). Both fields are used as input to our classifiers.

\section{Data Pre-Processing}
The analysis was restricted to English-language reviews. This follows the scope of the dataset specified for this study, which contains 1000 genuine and 1000 AI-generated English reviews. Restricting the analysis to English also allows the same linguistic preprocessing methods, such as English stop-word removal and WordNet-based lemmatization, to be applied consistently across all reviews. 

In the data, each review had both a positive and negative component. To get access to the full review we manually combined the negative and positive parts together. 

We divided the data into training and test sets using a 75\%--25\% split, where the vocabulary is built only from the training set, so no test information leaks into the features. Stratified sampling was used to ensure that the proportion of genuine and AI-generated reviews remained the same in both sets. A fixed random seed of 42 was used to make the split reproducible, so that the same reviews are assigned to the training and test sets each time the analysis is performed.

Before creating the bag-of-words, the texts were lower-cased. This was done to improve the data-set as e.g. the words "Hotel" and "hotel" should be treated as the same as the capitalization doesn't effect word function. The words were subsequently lemmatized using the NLTK wordnet lemmatizer. We did this because lemmatization often provides a more trainable dataset as words that have very close related meanings but differ in inflection. Making sure these words are identified the same reduces the number of features and combines related word forms. Stop words were hereafter removed to remove noise so more important terms wouldn't be drowned out.

The reviews were represented using a bag-of-words representation based on word frequencies. Each review is represented by the number of times each word occurs in that review, resulting in a sparse document-term matrix. Terms occurring in fewer than two training documents were removed using a minimum document frequency of two. Such very rare terms provide limited evidence for general classification because they occur in only a single training document. Removing these sparse terms therefore reduces the dimensionality of the feature space and can reduce the influence of features that are unlikely to generalize to unseen reviews.

\section{methods/setup}
To test the effectiveness of traditional machine learning methods in distinguishing between AI-generated reviews and genuine reviews, we employed 5 different models. Almost all models are unique in their classification, where we use a generative linear classifier, a discriminative linear classifier, a non-linear classifier, and two bagged ensemble tree models. In this way, we can test a variety of models and examine which one achieves the best performance metrics across accuracy, precision, recall, and F1 score. 

The first two models are linear models: Multinomial naive Bayes(MNB) and logistic regression. While both are linear classifiers, MNB is generative while logistic regression is discriminative. This difference means that while MNB predicts under the assumption that features are conditionally independent given the class, logistic regression directly models $P(class|features)$ and chooses weights by maximizing conditional likelihood. These models were tuned based on their smoothing $\alpha$ and $C$ values. 

The other three models are tree-based models. One single classifier tree, and two ensemble methods: random forests and gradient boosting. In comparison to the first two models, these three make no linear assumptions. They can handle nonlinearity and interactions simultaneously. The single classification being only one tree is the most sensitive to the data. Small changes can lead to completely different trees. The random forests on the other hand are trained on bootstrapped samples, and are averaged at the end, reducing bias significantly compared to the single tree. Gradient boosting builds shallow trees, where each new tree fits its errors of the current ensemble, which effectively reduces the bias. These models had more hyperparameters tuned versus the linear models, including: maximum depth, maximum features, and learning rate. 


\section{Results}

Throughout this section, AI-generated reviews are the positive class: precision
is the fraction of flagged reviews that really are AI-generated, and recall is
the fraction of AI-generated reviews that are detected.

\subsection{Hyperparameter Selection}

Hyperparameters were selected using 5-fold cross-validation and out-of-bag accuracy, where the latter was used only for the random forest model. Using these validation methods, we used only the 1500 training reviews, ensuring the 500 test samples would still be unseen for the model. The reason that the random forest method uses a different evaluation method is because of its use of bootstrapping. Each tree is grown on a bootstrap sample, so about a third (roughly 36.8\%) of the training reviews are not seen by a given tree. Predictions are then made only by the trees that did not see those specific training reviews, to give an estimate of the accuracy on new data. After the training, each model was then refitted on the full training set with the selected hyperparameters.

\begin{table}[H]
\centering
\small
\caption{Selected hyperparameters and validation accuracy on the training set.}
\label{tab:hyperparams}
\begin{tabular}{lll}
\hline
Model & Selected hyperparameters & Validation acc.\\
\hline
Naive Bayes & smoothing $\alpha = 0.1$ & 0.907 \\
Logistic regression & $C = 1$ & 0.895 \\
Single classification tree & max.\ depth $= 20$ & 0.798 \\
Random forest & 500 trees, max\_features = log2, min.\ split 20, min.\ leaf 1 & 0.916 \\
Gradient boosting & learning rate 0.1, max.\ depth 6, min.\ split 20, min.\ leaf 2 & 0.923 \\
\hline
\end{tabular}
\end{table}

\subsection{Test Set Performance}

Table~\ref{tab:performance} shows the results on the 500 test reviews. Naive Bayes has the highest accuracy (0.918) and F1 score (0.921), but Logistic Regression, the Random forest and Gradient Boosting are all within about two percentage points. The outlier model is the single tree, clearly underperforming relative to the other 4 models. Validation scores are close to the test score (e.g.\ 0.916 versus 0.914 test for the random forest), so the tuning did not seem to overfit. 

\begin{table}[H]
\centering
\small
\caption{Test set performance (500 reviews). Positive class: AI-generated.}
\label{tab:performance}
\begin{tabular}{lcccc}
\hline
Model & Accuracy & Precision & Recall & F1 \\
\hline
Naive Bayes         & \textbf{0.918} & 0.888 & \textbf{0.956} & \textbf{0.921} \\
Logistic regression & 0.910 & 0.929 & 0.888 & 0.908 \\
Classification tree & 0.794 & 0.808 & 0.772 & 0.789 \\
Random forest       & 0.914 & 0.929 & 0.896 & 0.912 \\
Gradient boosting   & 0.908 & \textbf{0.974} & 0.888 & 0.903 \\
\hline
\end{tabular}
\end{table}

\subsection{McNemar's Test}

Since all classifiers are evaluated on the same test reviews, we compare their accuracy with McNemar's test. For classifiers $A$ and $B$, let $n_{11}$ be the number of reviews both classify correctly, $n_{10}$ the number only $A$ gets right, $n_{01}$ the number only $B$ gets right, and $n_{00}$ the number both get wrong. Only the reviews on which they disagree are informative: if both have the same error rate, $n_{10}$ and $n_{01}$ should be about equal. The test statistic is
\[
\chi^2 = \frac{(|n_{10} - n_{01}| - 1)^2}{n_{10} + n_{01}},
\]
which approximately follows a chi-squared distribution with one degree of freedom under the null hypothesis. We use $\alpha = 0.01$. Table~\ref{tab:mcnemar} gives the $2 \times 2$ tables of all tests we performed.

\begin{table}[H]
\centering
\small
\caption{McNemar tests on the test set. Each row gives the $2\times2$ table
($n_{11}$: both correct, $n_{10}$: only $A$ correct, $n_{01}$: only $B$
correct, $n_{00}$: both wrong). NB = naive Bayes, LR = logistic regression,
CT = classification tree, RF = random forest, GB = gradient boosting.}
\label{tab:mcnemar}
\begin{tabular}{llrrrrrl}
\hline
$A$ & $B$ & $n_{11}$ & $n_{10}$ & $n_{01}$ & $n_{00}$ & $\chi^2$ & $p$ \\
\hline

NB & LR & 425 & 34 & 30 & 11 & 0.14 & 0.708 \\
RF & CT & 380 & 77 & 17 & 26 & 37.0 & $1.2 \times 10^{-9}$ \\
GB & CT & 387 & 79 & 10 & 24 & 52.0 & $5.7 \times 10^{-13}$ \\
RF & GB & 447 & 10 & 19 & 24 & 2.21 & 0.137 \\
GB & LR & 446 & 20 &  9 & 25 & 3.45 & 0.063 \\
\hline
\end{tabular}
\end{table}

\subsection{Generative versus Discriminative Linear Model}

Naive Bayes is slightly more accurate than logistic regression (0.918 versus 0.910), but the difference is not significant ($p = 0.708$): of the 64 review on which they disagree, Naive Bayes gets 34 right and logistic regression 30. The two models do make different kinds of errors. Naive Bayes labels about 270 test reviews as AI-generated, giving a high recall (0.956) but a lower precision (0.888). Logistic regression labels only about 240 reviews as AI-generated, with higher precision (0.929) and lower recall (0.888). Their F1 scores are similar, which then leads to the question which is truly preferable. The preference depends on whether falsely labeling a real review is worse than missing an AI-generated review.

\subsection{Single Tree versus Tree Ensembles}

Both tree ensembles perform significantly better than the single tree. Both models improve on all metrics, with the gradient boosting model improving the accuracy from a low 0.794 to 0.908 and the random forest model improving the accuracy even more to 0.914. Both $p$ values show significance in these results with $p$ values of $1.2 \times 10^{-9}$ and $5.7 \times 10^{-13}$ respectively. These results were expected with the bag-of-words data. A singular deep tree can only use a few of the many sparse terms while also having high variance. The random forest improves on this methodology by averaging many trees, while the gradient boosting model fits trees to the errors of the earlier ones. Averaging many trees reduces variance in the random forest, and boosting corrects the errors of earlier trees.. The $p$ value of the McNemar test also shows that both ensembles don't differ significantly from each other ($p=0.719$).

It is worth noting that while the ensemble methods do clearly outperform the single tree, when comparing the accuracy to the other two linear models, the accuracies do not differ too much, at most differing by only $0.01$. 

\subsection{Random Forest: Positive versus Negative Reviews}

The test set of 500 reviews contains 246 positive reviews and 254 negative reviews. The random forest model classifies 220 of the positive reviews correctly (accuracy 0.894), and 237 of the negative reviews correctly (accuracy 0.933). Since the positive and negative reviews are two separate groups, we used a chi-squared test of independence on the $2 \times 2$ table of sentiment versus correct/incorrect. This resulted in $\chi^2 = 1.92$ and $p = 0.166$. This shows us that model is not significantly better at classifying one sentiment over the other. Perhaps if the test set would have been significantly larger, then the model might have shown a clear difference. 

\subsection{Most Important Features}

For logistic regression we take the terms with the largest positive and
negative coefficients, and for naive Bayes the largest and smallest values of
$\log \hat{P}(w \mid \text{AI}) - \log \hat{P}(w \mid \text{human})$. For the
tree models we use the impurity-based importance, which has no direction, so we
checked in which class each term is more frequent (Table~\ref{tab:features}).

\begin{table}[H]
\centering
\small
\caption{Top-5 features per model. All tree-model terms are more frequent in
AI-generated reviews, except the one marked $+$.}
\label{tab:features}
\begin{tabular}{lll}
\hline
Model & Towards AI-generated & Towards human-written \\
\hline
Naive Bayes         & lacked, notch, desired,       & told, booked, said, \\
                    & inconsistent, bucharest       & com, 10 \\
Logistic regression & heart, limited, proximity,    & booked, construction, cleaning, \\
                    & lack, properly                & fan, big \\
\hline
Classification tree & \multicolumn{2}{l}{service, staff, hotel, noise, good$^+$} \\
Random forest       & \multicolumn{2}{l}{service, room, staff, hotel, attraction} \\
Gradient boosting   & \multicolumn{2}{l}{staff, convenient, located, service, attraction} \\
\hline
\end{tabular}
\end{table}

For the linear models, the AI-generated reviews tended to have stock phrases, such as ''left much to be \emph{desired}'', ''top-notch'', and would have generic complaints about a \emph{lack} of something. Another notable inclusion is the city name \emph{bucharest}. This is likely due to the generation prompt of the reviews specified the city name, whereas it would be very unlikely that a guest writing on a hotel's own Booking.com page would include the name of the city in the review. Real human-written reviews tended to have more clarity as to what happened: what a member of staff \emph{told} them or \emph{said}, or specific problems such as \emph{construction}, \emph{cleaning}, or a \emph{fan}. Both linear models agree that \emph{booked} points to a human written review, and \emph{lack}/\emph{lacked} points to AI-generated. 

While the linear models returned terms that are more correlated with one class or another, the nature of the tree models returns nothing directly about the direction of terms. The tree models, however, all returned more common words  such as \emph{service}, \emph{hotel}, \emph{room}, and \emph{staff}. This is likely due to GPT-4, which has gone through the standard aspects of a hotel stay, resulting in more common terms. The most important features for the tree models are also heavily skewed towards artificially generated reviews, as they may be longer on average than the real reviews. This results in the common words appearing more often, simply because there is more text in the fake reviews. 


\section{Discussion and Conclusion}

In this study we we compared five classifiers for detecting AI-generated hotel reviews only with word counts as our features. Even without any other over-engineered features, four of the five models classified more than 90\% of the test reviews correctly. This shows that the words used in the fake reviews are clearly distinguishable from the real ones in the words used. 

\paragraph{Generative versus discriminative.} The Naive Bayes and Logistic Regression models had almost the same accuracy (0.918 and 0.910), and the difference was not significant ($p = 0.708$). The models did differ in the type of error they made: Naive Bayes found more AI-generated reviews (recall 0.956 vs 0.888) but also flagged more genuine ones (precision 0.888 vs 0.929). The Logistic Regression model was clearly a bit more cautious, but which of the two models is more preferable depends on the cost of each kind of error.
 
\paragraph{Single tree versus ensembles.} 
Both tree ensembles were significantly more accurate than the single classification tree, which only reached an accuracy of 0.794. The random forest (0.914) and the gradient boosting (0.908) models did not differ significantly from each other ($p = 0.719$), nor in performance from the other two linear models. Clearly combining many trees outperforms a singular tree. But for this simple bag-of-words data, the extra complexity gave no significant gain over a simple linear model, showing that simplicity still can be an effective method.
 
\paragraph{Positive versus negative reviews.} 
The random forest model was slightly more accurate on negative reviews than on positive reviews (0.933 vs 0.894), but with $p = 0.166$, the results were not significant to show any real difference in classification capability. We therefore found no evidence that positive or negative reviews are easier to be classified. 


\paragraph{Important features.} 
The linear models provided more clarity in terms of the differences of the two classes. AI-generated reviews were more polished and generic (\emph{desired}, \emph{inconsistent}, \emph{limited}, \emph{proximity}, \emph{lack}, while genuine human-written reviews contained more specific events and problems (\emph{booked}, \emph{told}, \emph{said}, \emph{construction}, \emph{fan}). On the other hand, the tree models relied more on common words that would be commonly associated with a hotel stay such as \emph{service}, \emph{hotel}, \emph{room} and \emph{staff}, which could be due to how GPT-4 is tuned.

\paragraph{Limitations.} 
Our analysis has several limitations that should be addressed. Firstly, we used a single train/test split, which in combination with only 500 test reviews makes differences in accuracy almost negligible when using McNemar's test. Repeating the analysis over a larger sample of reviews and multiple random splits would likely give a better picture of the differences in results. 
Next, not every hyperparameter was tuned. Mainly the number of trees in the gradient boosting model was not tested through cross-validation due to hardware constraints.  
Thirdly, due to the nature of our preprocessing methodology, we did not include any n-grams, meaning that phrases such as ''left much to be desired'' would only be captured indirectly. For example the inclusion of bigrams would likely increase all metrics for all models. 
Lastly, this dataset was created by one model with one single prompt template. City names like \emph{bucharest} likely appeared because of how the generation prompt was structured. Using multiple models with various prompts would possibly make it more difficult for the models to classify AI-generated reviews, since each model has its own generation outputs. 

In conclusion, simple bag-of-words classifiers can detect GPT-4-generated hotel
reviews in this dataset with over 90\% accuracy. The linear models perform as
well as the more complex tree ensembles, and they are also easier to interpret.

\begin{thebibliography}{9}

\bibitem{ignat-etal-2025-maide}
Oana Ignat, Xiaomeng Xu, and Rada Mihalcea.
\newblock MAiDE-up: Multilingual Deception Detection of AI-generated Hotel Reviews.
\newblock \textit{Findings of the Association for Computational Linguistics: NAACL 2025}, pages 1636--1653, 2025.


\end{thebibliography}

\section{AI Statement}

Laurens: 
I used Claude to polish up some parts of the paper, mostly to ensure clarity and to check for any misspellings or typos. I also used Claude to generate an easier visualization in my code for the McNemar's test. The test already worked, but Claude made an easier table in the terminal output. It also helped explain how to implement out-of-bag evaluation.

\end{document}